% Manual do usuário — PCG Bids and Contracts (BOR2)
% Compilar com: pdflatex -interaction=nonstopmode pcg-bids-and-contracts-pt.tex
\documentclass[11pt,a4paper]{article}

\usepackage[T1]{fontenc}
\usepackage[utf8]{inputenc}
\usepackage[brazil]{babel}
\usepackage[margin=2.4cm,top=2.6cm,bottom=2.4cm]{geometry}
\usepackage{xcolor}
\usepackage{titlesec}
\usepackage{booktabs}
\usepackage{tabularx}
\usepackage{enumitem}
\usepackage{fancyhdr}
\usepackage{graphicx}
\usepackage{longtable}
\usepackage[hidelinks]{hyperref}

\definecolor{pcgblue}{HTML}{1D4ED8}
\definecolor{pcggray}{HTML}{4B5563}
\definecolor{pcglight}{HTML}{F3F4F6}
\definecolor{pcgamber}{HTML}{B45309}
\definecolor{pcggreen}{HTML}{15803D}

\hypersetup{
  colorlinks=false,
  pdftitle={Manual — PCG Bids and Contracts},
  pdfauthor={Premium Group},
}

\titleformat{\section}{\Large\bfseries\color{pcgblue}}{\thesection.}{0.6em}{}
\titleformat{\subsection}{\large\bfseries}{\thesubsection}{0.6em}{}
\titleformat{\subsubsection}{\normalsize\bfseries\color{pcggray}}{\thesubsubsection}{0.6em}{}
\titlespacing*{\section}{0pt}{1.6em}{0.8em}

\pagestyle{fancy}
\fancyhf{}
\fancyhead[L]{\small\color{pcggray} PCG Bids and Contracts — Manual do usuário}
\fancyhead[R]{\small\color{pcggray}\thepage}
\renewcommand{\headrulewidth}{0.4pt}

\setlist[itemize]{itemsep=2pt,topsep=4pt,leftmargin=1.4em}
\setlist[enumerate]{itemsep=3pt,topsep=4pt,leftmargin=1.6em}

% Rótulo de botão/campo da tela (a interface é em inglês)
\newcommand{\ui}[1]{\textbf{\textsf{#1}}}
% Caixa de aviso

\newsavebox{\avisobox}
\newenvironment{aviso}
  {\par\vspace{0.7em}\noindent
   \begin{lrbox}{\avisobox}\begin{minipage}{\dimexpr\linewidth-2.4em\relax}}
  {\end{minipage}\end{lrbox}%
   \noindent\colorbox{pcglight}{\hspace{0.6em}\usebox{\avisobox}\hspace{0.6em}}%
   \vspace{0.7em}\par}

\begin{document}

% ── Capa ──────────────────────────────────────────────────────────────────────
\begin{titlepage}
\centering
\vspace*{3cm}
{\Huge\bfseries\color{pcgblue} PCG Bids and Contracts\par}
\vspace{0.8cm}
{\LARGE Manual do usuário\par}
\vspace{0.4cm}
{\large Do cadastro do trade ao contrato assinado\par}
\vfill
{\large Premium Group — Business Operation Review\par}
{\color{pcggray} Versão 1.0 — \today\par}
\end{titlepage}

\tableofcontents
\newpage

% ── 1 ─────────────────────────────────────────────────────────────────────────
\section{O que é esta tela}

A tela \ui{PCG Bids and Contracts} é onde a PCG monta e acompanha a papelada de
cada serviço contratado numa obra. Ela produz \textbf{dois documentos}:

\begin{enumerate}
  \item \textbf{Bid Request} (pedido de cotação) — o formulário que vai para o
  subcontratado com a descrição do que se quer, para que ele devolva um preço.
  \item \textbf{Subcontract} (contrato) — o documento final, numerado, que o
  subcontratado assina.
\end{enumerate}

Entre um e outro existe um caminho: mandar, receber preço, aprovar, contratar,
assinar. A tela guarda cada passo desse caminho como um \textbf{evento}, com
data e autor. Nada é digitado duas vezes: o preço, o subcontratado e a situação
atual são \emph{lidos} desse histórico.

\subsection{Por que existe}

Antes, cada cotação era um arquivo do Word copiado do anterior. Isso significava
que o texto padrão variava de obra para obra, ninguém sabia dizer em que pé
estava cada serviço, e o contrato assinado não tinha como ser reproduzido depois
— o arquivo original já tinha sido editado por cima.

Aqui, o texto padrão mora num só lugar (o catálogo), e cada documento gerado
guarda uma cópia congelada do que foi usado. O contrato que o sub assinou em
agosto continua legível em outubro, mesmo que o padrão da empresa tenha mudado
no meio.

\subsection{Quem pode fazer o quê}

O acesso é dado em \ui{Settings > Users}, na permissão \ui{PCG Bids and Contracts}:

\begin{itemize}
  \item \textbf{Leitura} — vê tudo: projetos, catálogo e os documentos. Não
  altera nada.
  \item \textbf{Escrita} — cria projeto, responde questionário, registra evento
  e edita o catálogo.
\end{itemize}

Sem permissão de escrita, os botões de criar e editar simplesmente não aparecem.
Não é erro da tela.

% ── 2 ─────────────────────────────────────────────────────────────────────────
\section{Vocabulário}

Sete palavras aparecem o tempo todo. Vale ler esta página antes de mexer.

\begin{longtable}{@{}p{3.4cm}p{11.4cm}@{}}
\toprule
\textbf{Palavra} & \textbf{O que significa aqui} \\
\midrule
\endhead
\textbf{Projeto} & Uma obra. Tem nome, endereço, situação e tipo. Ex.: ``Bates
Quarry Lot 19''. \\[4pt]

\textbf{Trade} & Uma especialidade contratada: encanamento, elétrica, HVAC. Cada
trade dentro do projeto tem sua própria trilha, do começo ao fim, independente
das outras. \\[4pt]

\textbf{Catálogo} & O cadastro dos trades: as perguntas que serão feitas e o
texto padrão que entra nos documentos. É o molde; o projeto é a peça. \\[4pt]

\textbf{Questionário} & As perguntas do trade, respondidas para aquela obra
específica. É o que descreve o serviço. \\[4pt]

\textbf{Bid Request} & O formulário impresso a partir do questionário, enviado ao
sub para ele precificar. \\[4pt]

\textbf{Evento} & Um fato registrado: ``mandei'', ``recebi o preço'', ``aprovei''.
Cada evento tem data, autor e, quando faz sentido, valor e subcontratado. \\[4pt]

\textbf{Contrato} & O documento final, com número próprio, gerado quando um
preço foi aprovado. \\
\bottomrule
\end{longtable}

\begin{aviso}
\textbf{A regra de ouro:} a situação de um serviço nunca é escolhida à mão. Ela
é consequência do último evento registrado. Se algo aparece como ``Bid sent'',
é porque existe um evento de envio no histórico.
\end{aviso}

% ── 3 ─────────────────────────────────────────────────────────────────────────
\section{Antes de começar: o que precisa existir}

Esta é a parte que costuma ser pulada — e é justamente ela que faz o resto
funcionar. Nada aqui se repete por obra: é configurado uma vez e usado por todos
os projetos.

\subsection{O catálogo de trades}

Acesse \ui{PCG Bids and Contracts} e clique em \ui{Trades}, no topo da tela.
Você chega ao \ui{Trade Catalog}. Cada cartão é um trade cadastrado.

Para criar, clique em \ui{New trade}. Os campos:

\begin{itemize}
  \item \ui{Name} — o nome do serviço, como aparece em tudo. Ex.: ``Plumbing''.

  \item \ui{Document code} — três letras que identificam o trade no número do
  contrato. Ex.: \texttt{PLB} para Plumbing, \texttt{ELC} para Electrical. É
  daqui que sai o \texttt{PLB-00001} impresso no documento.

  \item \textbf{Ícone} — escolhido no botão ao lado do nome. Serve só para
  reconhecer o trade de relance na lista.

  \item \ui{Requires a bid form} — a chave mais importante da tela.
  \begin{itemize}
    \item \textbf{Ligada:} o serviço passa por cotação. Existe formulário, envio
    ao sub, recebimento de preço e aprovação, e só então o contrato.
    \item \textbf{Desligada:} não há cotação. O contrato é gerado direto do
    escopo padrão. Use para serviço de preço já acordado ou de rotina.
  \end{itemize}

  \item \ui{Standard note} — um texto que vale para todo projeto desse trade.
  Aparece no formulário de cotação, acima das perguntas.
\end{itemize}

\subsubsection{As perguntas}

Dentro do editor do trade, abra a seção \ui{Questions}. Cada pergunta tem:

\begin{itemize}
  \item \textbf{Enunciado} — o que se pergunta. Ex.: ``Tipo de aquecedor''.
  \item \textbf{Tipo} — quatro opções:
  \begin{itemize}
    \item \ui{Single choice} — escolhe uma entre várias.
    \item \ui{Yes / No} — sim, não, ou ``a definir''.
    \item \ui{Multiple choice} — pode marcar mais de uma.
    \item \ui{Free text} — texto livre.
  \end{itemize}
  \item \textbf{Opções} — a lista de respostas possíveis (não se aplica a texto livre).
  \item \ui{Hint} — uma orientação curta que aparece embaixo do campo.
  \item \textbf{Etiqueta} — \ui{All projects} (respondida em toda obra) ou
  \ui{Optional} (só quando o projeto pedir).
  \item \textbf{Quantidade} — quando ligada, a pergunta ganha um segundo campo
  numérico, usado só se a resposta for afirmativa. O \ui{Hint} vira a unidade
  dessa contagem: ``cômodos'', ``dutos''.
\end{itemize}

\subsubsection{O escopo padrão}

Ainda no editor, a seção de conteúdo do documento traz três listas. Elas são
copiadas para todo documento gerado desse trade:

\begin{itemize}
  \item \ui{Work included} — o que está incluído no serviço.
  \item \ui{Exclusions} — o que explicitamente \emph{não} é responsabilidade do sub.
  \item \ui{Responsibility matrix} — quem fornece o quê: a construtora, o sub, ou
  ``a definir na assinatura''.
\end{itemize}

\subsubsection{As regras condicionais}

Esta é a parte que dá poder ao catálogo. Uma \textbf{regra} liga uma resposta do
questionário a uma mudança no escopo do contrato. Assim o mesmo trade produz
documentos diferentes conforme a obra, sem ninguém reescrever nada à mão.

Cada regra se lê como uma frase: \emph{``Se a pergunta X for Y, então faça Z na
lista W''}.

\begin{longtable}{@{}p{3.6cm}p{11.2cm}@{}}
\toprule
\textbf{Parte} & \textbf{Opções} \\
\midrule
\endhead
Pergunta & Qualquer pergunta do próprio trade. \\[3pt]
Condição & \ui{is} (é exatamente), \ui{includes} (contém, para múltipla escolha),
\ui{is answered} (foi respondida), \ui{is not answered} (ficou em branco). \\[3pt]
Ação & \ui{Add clauses} (acrescenta cláusulas), \ui{Remove a clause} (tira uma
cláusula do padrão), \ui{Replace a clause} (troca uma pela outra). \\[3pt]
Lista & \ui{Work included}, \ui{Exclusions} ou \ui{Responsibility matrix}. \\
\bottomrule
\end{longtable}

Uma regra pode acrescentar várias cláusulas de uma vez — a relação não é de uma
para uma.

\begin{aviso}
\textbf{Exemplo prático.} Se a pergunta ``O sub fornece o material?'' for
respondida com ``Não'', uma regra pode remover a cláusula padrão que diz que o
fornecimento é dele, e acrescentar em \ui{Exclusions} a linha dizendo que o
material será fornecido pela construtora. Quem responder o questionário não
precisa saber que isso existe: acontece sozinho.
\end{aviso}

\subsection{Os módulos do documento}

De volta ao \ui{Trade Catalog}, o botão \ui{Document defaults} abre o miolo dos
documentos — os blocos de texto que valem para todos os trades. É aqui que moram
as condições gerais, as cláusulas de seguro, os prazos de pagamento padrão.

Cada módulo tem:

\begin{itemize}
  \item \textbf{Título} e \textbf{corpo} — o texto em si.
  \item \textbf{Onde aparece} — \ui{Bid requests}, \ui{Contracts} ou \ui{Both}.
  \item \textbf{Tipo}:
  \begin{itemize}
    \item \ui{Editable} — texto que você escreve e pode alterar.
    \item \ui{Generated} — desenhado pelo sistema a partir dos dados. São três:
    \emph{Exhibit A: Scope of Work} (o escopo montado das respostas e das
    regras), \emph{Payment schedule} (o cronograma de pagamento) e
    \emph{Signatures} (os campos de assinatura).
    \item \ui{Standard} — texto fixo, que ninguém edita.
  \end{itemize}
  \item \textbf{Numeração} — se o módulo recebe número de seção. Apagar a seção 4
  renumera as seguintes; não fica buraco.
  \item \textbf{Ordem} — a ordem da lista é a ordem impressa. Arraste para mudar.
\end{itemize}

O cabeçalho com o logo, a tabela de dados da obra e o rodapé não são módulos: são
a moldura do papel e não se editam.

\subsection{O cadastro do subcontratado}

A lista de subs \textbf{não} é mantida aqui. Ela vem da tela
\ui{Subcontractor Docs}. Um sub cadastrado lá aparece aqui automaticamente, e
subs marcados como inativos não aparecem.

O documento precisa de três dados além do nome da empresa: \textbf{responsável},
\textbf{e-mail} e \textbf{telefone}. Quando falta algum, o botão ao lado do nome
do sub fica âmbar — é o aviso de que o contrato vai sair com um campo vazio.
Clique nele para completar.

% ── 4 ─────────────────────────────────────────────────────────────────────────
\section{Passo a passo: do zero ao contrato assinado}

Com o catálogo pronto, o trabalho do dia a dia são nove passos.

\subsection{Passo 1 — Criar o projeto}

Na tela principal, clique em \ui{New project}. O projeto nasce vazio, com o nome
``Untitled project'', e o painel de edição abre sozinho. \textbf{Não existe botão
de salvar}: cada campo é gravado assim que você sai dele.

Preencha:

\begin{itemize}
  \item \ui{Title} — o nome da obra. Ex.: ``Bates Quarry Lot 19''.
  \item \ui{Address} — rua, cidade, estado e CEP.
  \item \ui{Status} — \ui{Active}, \ui{On hold} ou \ui{Completed}. Serve para
  filtrar a lista depois.
  \item \ui{Type} — \ui{New Construction}, \ui{Addition} ou \ui{Renovation}.
  Perguntado uma vez para a casa inteira, e impresso logo abaixo do endereço nos
  dois documentos.
\end{itemize}

\subsection{Passo 2 — Adicionar os trades}

No mesmo painel, na seção \ui{Trades on this project}, clique em \ui{Add trade} e
escolha os serviços que essa obra precisa. Cada trade adicionado vira uma trilha
independente: pode estar em cotação enquanto o vizinho já está assinado.

\subsection{Passo 3 — Responder o questionário}

Abra o trade e vá em \ui{Questionnaire}. O contador ao lado do título mostra
quantas perguntas já foram respondidas — \texttt{7/12}. Ele fica verde quando
todas estão.

Ao final há o campo \ui{Additional Notes}: um texto livre para o que as perguntas
não cobrem. É opcional, não conta para o contador, mas \textbf{é impresso no
documento} e faz parte do que o sub está precificando.

\begin{aviso}
\textbf{Enquanto o questionário não estiver completo, nada sai.} Os botões de
gerar documento ficam bloqueados com a mensagem \emph{``Answer all N questions
first''}. Isso é proposital: um formulário incompleto vira preço errado.
\end{aviso}

\subsection{Passo 4 — Gerar o Bid Request}

Com o questionário completo, clique em \ui{Bid request}. Abre uma pré-visualização
do documento exatamente como ele será impresso. O botão
\ui{Print / Save as PDF} manda para a impressora ou salva em PDF, em A4.

O Bid Request imprime todas as opções de cada pergunta, com as marcadas
assinaladas — é um formulário, e o sub precisa ver o que \emph{não} foi pedido
tanto quanto o que foi.

\subsection{Passo 5 — Registrar o envio ao sub}

Enviado o PDF ao subcontratado (por e-mail, por fora do sistema), volte ao trade,
abra o \ui{History} e clique em \ui{Log event}. Escolha \ui{Bid sent to sub},
selecione o subcontratado e confirme a data.

\textbf{O mesmo formulário pode ir para vários subs.} Registre um envio para cada
um. O conjunto de envios do mesmo formulário é chamado de \textbf{rodada}, e o
painel \ui{Bids out} mostra a rodada inteira: quem recebeu, quem já respondeu,
por quanto.

\subsection{Passo 6 — Registrar o preço recebido}

Quando o sub responde, registre \ui{Bid received from sub}. Dois campos:

\begin{itemize}
  \item \ui{Bid amount} — o preço que ele mandou.
  \item \ui{Term set by PCG} — o prazo, em dias, semanas ou meses. Atenção: esse
  prazo é \textbf{definido pela PCG}, não cotado pelo sub. É quanto tempo a
  empresa dá para o serviço.
\end{itemize}

\subsection{Passo 7 — Aprovar e montar o pagamento}

Com os preços na mão, escolha um. Registre \ui{Bid approved} no sub vencedor.
Nos demais, registre \ui{Bid declined}.

\textbf{Só um bid pode ser aprovado.} É a aprovação que diz contra qual preço o
contrato será escrito.

Junto da aprovação vem o \textbf{cronograma de pagamento}. Clique no painel do
cronograma e monte as parcelas:

\begin{itemize}
  \item Cada linha tem um marco (ex.: ``Rough-in completo'') e um
  \textbf{percentual}.
  \item Os valores em dinheiro são calculados sozinhos, a partir do preço
  aprovado. Você nunca digita dinheiro aqui — assim o cronograma e o preço não
  têm como discordar.
  \item O sistema mostra quanto falta para fechar 100\%. Enquanto não fechar, o
  contrato não sai.
\end{itemize}

\subsection{Passo 8 — Gerar o contrato}

Com bid aprovado e cronograma montado, o botão \ui{Contract} libera. Ao clicar
pela primeira vez, o sistema \textbf{carimba o número do contrato}: três letras
do trade e cinco dígitos de um contador único da empresa — \texttt{PLB-00001}.

\begin{aviso}
Esse número é gravado uma vez e nunca muda. Se outro contrato for criado,
apagado ou reordenado depois, o número do papel que alguém assinou continua o
mesmo. Os cinco dígitos dizem quantos contratos a PCG já emitiu no total; as
letras dizem de qual trade este é.
\end{aviso}

Diferente do Bid Request, o contrato imprime apenas o que foi \emph{escolhido} —
uma opção não marcada não é cláusula de acordo nenhum.

\subsection{Passo 9 — Enviar, ajustar e assinar}

\begin{enumerate}
  \item \ui{Contract sent to sub} — registre quando o contrato for enviado.

  \item \ui{Contract adjusted} — se o sub contestar alguma cláusula. Veja a
  seção~\ref{sec:ajuste} sobre como reescrever o texto só para esse contrato.
  \textbf{Pode ser registrado uma única vez}, e só antes da assinatura.

  \item \ui{Contract signed} — o fim da linha. Registre com o
  \ui{SharePoint link} do PDF assinado, para que qualquer pessoa consiga chegar
  ao documento original a partir daqui.
\end{enumerate}

% ── 5 ─────────────────────────────────────────────────────────────────────────
\section{Os eventos, um a um}

\begin{longtable}{@{}p{3.9cm}p{5.3cm}p{5.2cm}@{}}
\toprule
\textbf{Evento} & \textbf{Quando usar} & \textbf{O que ele pede} \\
\midrule
\endhead
\ui{Trade opened} & Automático, ao adicionar o trade. & Nada. \\[4pt]

\ui{Bid sent to sub} & O formulário foi enviado a um sub. Um por sub. &
Subcontratado e data. \\[4pt]

\ui{Bid received} & O sub devolveu o preço. & Valor e prazo definido pela PCG. \\[4pt]

\ui{Bid approved} & O preço foi aceito. Só um por trade. & Subcontratado e o
cronograma de pagamento. \\[4pt]

\ui{Bid declined} & O preço de um sub foi recusado. & Subcontratado. \\[4pt]

\ui{Bid adjusted by the sub} & O sub venceu, mas depois avisou que não faz parte
do serviço. Carrega a aprovação para as respostas novas, em vez de rasgá-la. &
Nota explicando. \\[4pt]

\ui{Contract sent to sub} & O contrato foi enviado. & Subcontratado e data. \\[4pt]

\ui{Contract adjusted} & O sub pediu mudança de cláusula. Uma vez só, antes da
assinatura. & Nota, e opcionalmente o novo cronograma. \\[4pt]

\ui{Contract signed} & Contrato assinado. & Link do SharePoint. \\
\bottomrule
\end{longtable}

\subsection{O que dá para corrigir depois}

De um evento já registrado, dá para acertar: a data, a nota, o valor, o prazo, o
link e o subcontratado.

Não dá para mudar: \textbf{o tipo do evento} (um bid enviado não vira outra coisa
retroativamente) e \textbf{quem registrou e quando} — essa é a trilha de
auditoria.

A data só se move dentro da janela que os eventos vizinhos permitem: um preço não
pode ter voltado antes de o formulário ter saído. O calendário já esconde os dias
inválidos.

% ── 6 ─────────────────────────────────────────────────────────────────────────
\section{Por que o botão está bloqueado}

Todo bloqueio da tela tem uma frase explicando. Esta é a tradução delas.

\begin{longtable}{@{}p{7.2cm}p{7.4cm}@{}}
\toprule
\textbf{A tela diz} & \textbf{Significa} \\
\midrule
\endhead
\emph{Answer all N questions first} & O questionário está incompleto. Falta
responder alguma pergunta. \\[4pt]

\emph{The bid has to be approved first} & Não existe aprovação. Registre
\ui{Bid approved} no sub escolhido. \\[4pt]

\emph{The answers changed — the bid needs a new approval} & Alguém alterou uma
resposta depois da aprovação. Veja a seção seguinte. \\[4pt]

\emph{Set the payment schedule on the approved bid first} & Falta montar o
cronograma, ou ele não fecha 100\%. \\[4pt]

\emph{Nothing to log at this point} & O passo seguinte depende de algo que ainda
não aconteceu. \\[4pt]

\emph{This trade is finished} & O contrato já está assinado. Não há mais passo. \\
\bottomrule
\end{longtable}

% ── 7 ─────────────────────────────────────────────────────────────────────────
\section{Situações comuns}

\subsection{Mudei uma resposta depois de o bid ter sido aprovado}

O sistema percebe. A aprovação vale para \emph{aquele} formulário; alterada a
resposta, o papel deixou de ser o mesmo, e o contrato trava com o aviso
\emph{``the bid needs a new approval''}.

Dois caminhos honestos:

\begin{itemize}
  \item \textbf{Mandar de novo.} A rodada recomeça: novo \ui{Bid sent}, novo
  preço, nova aprovação. É o certo quando a mudança altera o preço.
  \item \textbf{Registrar \ui{Bid adjusted by the sub}.} Serve quando quem pediu a
  mudança foi o próprio sub que venceu — ele avisou que não faz parte do serviço.
  A aprovação acompanha as respostas novas em vez de ser descartada.
\end{itemize}

\subsection{O sub quer mudar uma cláusula do contrato}
\label{sec:ajuste}

Um sub discutindo uma cláusula não é motivo para mudar o padrão da PCG. Use o
editor de termos do contrato: ele reescreve o texto \textbf{apenas deste
contrato}.

O que fica guardado é a diferença. Um módulo que ninguém tocou continua seguindo
o catálogo, mesmo que o catálogo seja reescrito depois. E o módulo alterado
continua com o texto negociado, mesmo assim.

Registre também o evento \ui{Contract adjusted}, para que o histórico mostre que
esse contrato foi renegociado.

\subsection{Vários subs cotando o mesmo serviço}

É o funcionamento normal. Registre um \ui{Bid sent} para cada, e um
\ui{Bid received} conforme os preços chegarem. O painel \ui{Bids out} monta a
tabela comparativa da rodada. No fim, um \ui{Bid approved} e os demais
\ui{Bid declined}.

\subsection{Serviço que não passa por cotação}

É o trade com \ui{Requires a bid form} desligado. A trilha encurta: responde o
questionário e gera o contrato direto. Não existe botão de Bid request, nem
envio, nem aprovação — porque não há o que aprovar.

\subsection{Apagar um projeto}

Some tudo junto: as respostas de todos os trades, os subs e o estado dos
contratos. A tela pede confirmação por isso. Se o contrato já foi assinado, o PDF
no SharePoint continua lá — mas o histórico daqui não volta.

% ── 8 ─────────────────────────────────────────────────────────────────────────
\section{Resumo de uma página}

\begin{enumerate}
  \item \textbf{Uma vez só:} cadastre os trades (perguntas, escopo, regras) e os
  módulos do documento.
  \item \textbf{Por obra:} crie o projeto e adicione os trades.
  \item \textbf{Por trade:} responda o questionário até o contador ficar verde.
  \item Gere e envie o \textbf{Bid Request}; registre \ui{Bid sent} para cada sub.
  \item Registre \ui{Bid received} com preço e prazo.
  \item \ui{Bid approved} no vencedor, \ui{Bid declined} nos outros, e monte o
  cronograma até fechar 100\%.
  \item Gere o \textbf{contrato} — o número é carimbado aqui.
  \item \ui{Contract sent}, ajuste se precisar, e \ui{Contract signed} com o link
  do arquivo assinado.
\end{enumerate}

\vfill
\noindent{\color{pcggray}\small
Dúvidas ou erro no documento: procure o time de Business Operation Review.}

\end{document}
